\begin{table}[ht]
\centering
\caption{\label{tab:descriptive-statistics}Descriptive statistics. \label{tab:desc}}
\centering
\begin{threeparttable}
\begin{tabular}[t]{llrrrr}
\toprule
\multicolumn{2}{c}{ } & \multicolumn{2}{c}{a. UK (N=3968)} & \multicolumn{2}{c}{b. Other EU (N=11456)} \\
\cmidrule(l{3pt}r{3pt}){3-4} \cmidrule(l{3pt}r{3pt}){5-6}
  &    & Mean & Std. Dev. & Mean & Std. Dev.\\
\midrule
Number of affiliates in EU &  & 5.1 & 6.6 & 5.2 & 6.2\\
Affiliate size &  & 164.0 & 695.5 & 163.8 & 625.6\\
Japanese ownership ratio &  & 0.7 & 0.4 & 0.7 & 0.4\\
Year &  & 2015.8 & 3.3 & 2015.8 & 3.3\\
\midrule\\
 &  & N & Pct. & N & Pct.\\
High-income EU & High-income EU & 3968 & 100.0 & 9584 & 83.7\\
 & Other EU & 0 & 0.0 & 1872 & 16.3\\
Exit dummy & Exit & 1052 & 26.5 & 3001 & 26.2\\
 & Survive & 2916 & 73.5 & 8455 & 73.8\\
Sector & Agriculture \& Mining & 88 & 2.2 & 96 & 0.8\\
 & Manufacturing & 872 & 22.0 & 3080 & 26.9\\
 & Services & 3008 & 75.8 & 8280 & 72.3\\
\bottomrule
\end{tabular}
\begin{tablenotes}
\item \textit{Note: } 
\item This table shows the descriptive statistics of the estimation sample. The sample includes the Japanese affiliates in the EU. The variable `Number of affiliates in EU` is the number of affiliates in the EU.The variable `Affiliate size` is the size of the affiliate. The variable `High-income EU` is a dummy variable that takes 1 if the affiliate is located in the high-income EU countries. The variable `Exit` is a dummy variable that takes 1 if the affiliate exits the host country. The variable `Sector` is the sector of the affiliate.
\end{tablenotes}
\end{threeparttable}
\end{table}
